\section{Courbes d'évolution des paramètres étudiés durant l'expérimentation}
\subsection{Résultats des paramètres climatiques et hydriques et leurs relations}

\begin{figure}[H]
    \centering
    \imgstd{images/f_evolution_chap5.jpg}
    \caption{Courbes d'évolution de la température, de l’humidité, du vent et de ET0 en fonction des jours. }
    \label{fig:H_T_vent_}
\end{figure}

La température maximale moyenne se situe autour de 34,5 °C(30--39 °C, maximum observé le 10 août), tandis que la température minimale reste constante entre 25 et 27 °C, affichant une variation quotidienne moyenne de 8,6 °C. En moyenne, l'humidité relative fluctue de 57,6~\% au minimum à 91,5~\% au maximum, enregistrant jusqu'à 10 jours où elle atteint ou dépasse 95~\%. La vitesse moyenne du vent atteint 12~km/h, avec un pic enregistré à 21,3~km/h le 30 juillet. L'évapotranspiration moyenne s'élève à 5,28~mm/j, atteignant un minimum le 26/07 (avec 3,68~mm/j) et un maximum le 03/08 (avec 7,18~mm/j). Malgré des périodes humides notables, $ET_0$ varie relativement peu (5,6~mm/j pour fin juillet, 5,1 pour août et 5,2 pour septembre).

\begin{figure}[H]
    \centering
    \imgstd{images/g_kc_chap5.png}
    \caption{Evolution du Kc moyen selon la culture. }
    \label{fig:Kc_moyen5}
\end{figure}
Le navet a le Kc moyen le plus élevé (0,98) car son stade mi-saison (Kc = 1,0) est atteint plus tôt (J26) et dure plus longtemps. Le gombo et la laitue ont un Kc moyen de 0,89.

\begin{figure}[H]
    \centering
    \imgstd{images/h_ET0_ETc_chap5.png}
    \caption{Comparaison entre l'évolution de ET0 et de ETc selon la culture. }
    \label{fig:eto-etc-moyen5}
\end{figure}
 L’$ET_0$ est commune aux trois cultures (5,3 à 5,5~mm/j), puisque le climat est le même. L'ETc varie selon le $K_c$ : 4,74 (gombo), 4,61 (laitue) et 5,39~mm/j (navet). Le besoin en eau du navet est donc le plus proche de la demande climatique ($ET_c$/$ET_0$ de 0,98, contre 0,88 à 0,90 pour les deux autres).

\begin{figure}[H]
    \centering
    \imgstd{images/i_evapotranspiration_chap5.jpeg}
    \caption{Relation entre $ET_0$, $ET_c$ et la dose d’irrigation réellement apportée. }
    \label{fig:eto-etc-Di5}
\end{figure}
Avec l'apparition de la saison des pluies, $ET_0$ descend de 9,5 à 4-5~mm/j, tandis que $ET_c$ s'aligne sur le $K_c$ de chaque type de culture. Suite au 17/07, les apports d’irrigation deviennent plus rares : quasi inexistantes pour le gombo, principalement en début de croissance pour la  laitue, et de plus en plus fréquentes pour le navet.

\begin{figure}[H]
    \centering
    \imgstd{images/j_cumule_Di_chap5.jpeg}
    \caption{Relation entre la dose d’irrigation journalière et la dose d’irrigation réellement appliquée. }
    \label{fig:Di-Di-aporter5}
\end{figure}
Sur l’ensemble du cycle cultural, les doses d’irrigation réellement appliquées représentent 72~\% de la dose théorique pour le gombo, 76~\% pour la laitue et 34~\% pour le navet. Ces proportions sont calculées en rapportant la dose d’irrigation réellement appliquée à la dose théorique, puis en multipliant le résultat par 100.

\begin{figure}[H]
    \centering
    \imgstd{images/k_humidite_chap5.png}
    \caption{Evolution de l’humidité moyenne du sol selon la culture. }
    \label{fig:evolution-hum}
\end{figure}
La laitue est la plus riche en humidité, suivie par le gombo puis le navet, avec respectivement 28,2~\%, 24,6~\% et 21,6~\% d'humidité. L’ordre est cohérent avec l'irrigation constatée : le navet est le plus réclamé et le plus aride. La laitue présente également la variance la plus importante (écart-type de 10~\%, pic à 55,9~\% le 20/08). Cette élévation observée au sein du plot test de la laitue s’explique par sa position topographique au niveau de la parcelle.

\begin{figure}[H]
    \centering
    \imgstd{images/l_ph_chap5.png}
    \caption{Evolution du pH moyen du sol selon la culture. }
    \label{fig:ph-sol5}
\end{figure}
Le pH du sol est neutre à légèrement alcalin: 7,55 pour le gombo, 7,89 pour le navet, et 8,23 pour la laitue. Les divergences (approximativement 0,7 unité) demeurent en dessous de la variabilité quotidienne(écart-type approximatif de 1 unité). L'axe débute à 6, ce qui amplifie visuellement les écarts.

\begin{figure}[H]
    \centering
    \imgstd{images/m_EC_chap5.png}
    \caption{Evolution de la EC  moyenne du sol selon la culture. }
    \label{fig:ec-moyensol5}
\end{figure}
Les résultats de l'ANOVA montrent que la différence entre les moyennes de conductivité électrique des trois cultures n'est pas statistiquement significative ($F(2,\,122) = 0{,}79$ ; $p = 0{,}46$). Les moyennes observées sont de 298,4~\textmu S/cm pour la laitue ($n = 42$), 276,3~\textmu S/cm pour le gombo ($n = 55$) et 260,6~\textmu S/cm pour le navet ($n = 28$). Malgré ces différences numériques, le test ne met pas en évidence de différence significative entre les cultures au seuil de 5~\%.

\begin{figure}[H]
    \centering
    \imgstd{images/n_NPK_chap5.png}
    \caption{Evolution moyenne des teneurs en N, P, K du sol selon la culture. }
    \label{fig:evolutionNPK5}
\end{figure}
Les teneurs estimées en potassium sont les plus élevées, suivies de celles du phosphore puis de l’azote. Pour le gombo ($n = 55$), les moyennes sont de 13,35~mg/kg pour N, 18,86~mg/kg pour P et 43,74~mg/kg pour K. Pour la laitue ($n = 42$), elles sont respectivement de 14,40, 20,39 et 47,26~mg/kg. Pour le navet ($n = 28$), elles sont de 12,55, 17,73 et 41,19~mg/kg. Ces valeurs étant dérivées de la EC, elles ne constituent pas des mesures indépendantes des teneurs en $N$, $P$ et $K$. Le graphique est donc interprété de manière descriptive.

\begin{figure}[H]
    \centering
    \imgstd{images/o_boxplot_ph_chap5.png}
    \caption{Evolution de la variabilité journalière du pH du sol selon la culture. }
    \label{fig:ph-boxplot5}
\end{figure}
Le $pH$ présente une variabilité journalière dans les trois cultures, avec des valeurs comprises entre 4{,}6 et 10{,}0. Les moyennes sont de 7{,}55 pour le gombo ($n = 55$), de 8{,}23 pour la laitue ($n = 42$) et de 7{,}89 pour le navet ($n = 28$). La laitue présente ainsi le pH moyen le plus élevé, tandis que le gombo présente le plus faible. Ces différences sont décrites à partir des observations disponibles, sans conclure à leur significativité statistique en l’absence d’un test comparatif.

\begin{figure}[H]
    \centering
    \imgstd{images/p_evolution_ph_EC_chap5.png}
    \caption{Evolution temporelle du pH et de la EC du sol pour la laitue. }
    \label{fig:ph-ec-evolution5}
\end{figure}
Au cours du cycle cultural de la laitue ($n = 42$ observations journalières), la conductivité électrique moyenne augmente de 190~\textmu S/cm en juillet à 293~\textmu S/cm en août, puis à 375~\textmu S/cm en septembre, avec un maximum journalier de 898~\textmu S/cm le 8 septembre. Le $pH$ passe de 7{,}3 en juillet à environ 8{,}4 en août et en septembre, tout en présentant des fluctuations journalières et deux pics à 10. L’évolution de l’$EC$ et du $pH$ est décrite à partir des observations recueillies. Aucune conclusion statistique sur leur relation avec l’humidité n’est formulée, faute de coefficient de corrélation et de valeur $p$ disponibles.

\begin{figure}[H]
    \centering
    \imgstd{images/q_tem_hum_ph_chap5.png}
    \caption{Evolution temporelle de l'humidité, de la température et du pH du sol selon la culture. }
    \label{fig:ph-hum-temp5}
\end{figure}
Les températures moyennes du sol restent globalement comprises entre 29 et 37~°C, avec peu de différences apparentes entre les cultures. Pour la laitue, l’humidité du sol augmente d’environ 18~\% en juillet à 34~\% en septembre, une évolution observée pendant la saison des pluies. Le $pH$ présente des fluctuations au cours du cycle, avec une hausse apparente pour la laitue. Ces tendances sont descriptives et ne permettent pas, à elles seules, d’établir un lien statistique entre les paramètres étudiés.

\begin{figure}[H]
    \centering
    \imgstd{images/r_relation_ph_EC_chap5.png}
    \caption{Nuage de point mettant en relation le pH et l’EC du sol. }
    \label{fig:ph-ec-nuage5}
\end{figure}
Le nuage de points permet d’examiner visuellement la relation entre le $pH$ et la $EC$ du sol. Les coefficients de corrélation descriptifs obtenus sont de +~0,40 pour la laitue, +~0,08 pour le gombo et -~0,22 pour le navet. Ces valeurs indiquent des associations de sens et d’intensité différents selon la culture. Toutefois, en l’absence des effectifs exacts utilisés pour chaque corrélation et des résultats des tests de significativité, il est impossible de conclure à l’existence d’une relation statistiquement significative entre l'$EC$ et le $pH$.

\begin{figure}[H]
    \centering
    \imgstd{images/s_ph_EC_duSol_chap5.png}
    \caption{Relation de corrélation entre les paramètres du sol et du climat. }
    \label{fig:correlation5}
\end{figure}
Il existe une parfaite corrélation entre $EC$, l'azote, le phosphore et le potassium ($r = 1$), ce qui les rend redondants. $ET_0$ est principalement définie par l'humidité relative minimale ($r = -0{,}91$), la température maximale ($r = +0{,}80$) et le vent ($r = +0{,}63$). La corrélation entre l'humidité du sol et $EC$ est positive ($r = +0{,}54$), tandis qu'avec $ET_0$, elle est négative ($r = -0{,}49$), et l'association avec la température est faible ($r = -0{,}22$).
Les corrélations sont évaluées sur les trois cultures mélangées, et celles entre $ET_c$ et le $K_c$ (0,70) reflètent principalement la progression du stade de croissance. L'évapotranspiration est le principal moteur derrière la dynamique de l'eau. Les amendements ont un impact direct sur la salinité et la fertilité. L'efficacité du contrôle de l'irrigation nécessite l'inclusion de ces corrélations.

\subsection{Consommation d’eau : zone test contre zone témoin}

La figure ci-dessous compare le volume d’eau apporté par plot de $2\,\mathrm{m^2}$ sur le cycle de chaque culture. Le pilotage par capteur a réduit la consommation d’eau de 86~\% pour le gombo (173 L contre 1 281 L), de 76~\% pour la laitue (99 L contre 420 L) et de 66~\% pour le navet (373 L contre 1 113 L). Sur la zone témoin, l’apport est régulier, de l’ordre de 21 litres par jour, avec des plateaux lors des périodes de pluie, tandis que sur la zone test la courbe ne progresse que lorsque le capteur détecte un besoin en eau.

Cet écart est cohérent avec le principe du pilotage, qui ajuste les apports à l’état d’humidité du sol plutôt qu’à un volume fixe. La consommation plus élevée du navet parmi les plots test pourrait être liée à son développement foliaire (transpiration), sans que cela ait été vérifié. Concernant la fertilisation, l’amendement ayant été réalisé avant les cultures, la fertilisation complémentaire n’a été nécessaire que lors des pluies les plus marquées et pendant la production ; ces apports n’ont pas été chiffrés.

Ces résultats valent pour les conditions de cet essai : le témoin de la laitue a été arrêté après 20 jours (plantes mortes), l’économie de 76~\% est donc un minimum, et l’économie d’eau peut aussi refléter la meilleure rétention du sol amendé, sans qu’on puisse séparer les deux effets.

\begin{figure}[H]
    \centering
    \imgstd{images/eau_appliquee_temoin_test_comparaison.jpeg}
    \caption{Comparaison de la consommation d'eau entre test et témoin. }
    \label{fig:eauTotal}
\end{figure}

\subsection{Résultats des paramètres de croissance et de développement végétative}
\begin{figure}[H]
    \centering
    \imgstd{images/t_survivant_chap5.jpeg}
    \caption{Evolution du nombre de plantes survivantes en fonction du jour. }
    \label{fig:survivant5}
\end{figure}
La laitue témoin disparaît complètement à J20, tandis que le test lui, garde ses plants à J20, puis 5 à J40. Le gombo témoin chute de 9 à 1 plant entre J14 et J27 soit 89~\% de mortalité, pendant que le test reste stable à 4 plants de J50 à J70. Le navet témoin passe de 6 à 3 plants entre J38 et J67 or le test reste stable.

\begin{figure}[H]
    \centering
    \imgstd{images/u_nbre_feuille_chap5.png}
    \caption{Evolution moyenne du nombre de feuille en fonction du jour. }
    \label{fig:nbre-feuille5}
\end{figure}
Le gombo test atteint en moyenne  9,75 feuilles à J70, contre 5 pour le témoin, soit une hausse de 95~\%. Le navet test monte à 14,8 feuilles à J50  contre 13,0 à J67 du côté témoin (mais avec seulement 3 plants survivants). La laitue test atteint 9,8 feuilles à J40.

La chute du nombre de feuilles chez le GOM TS s’explique par l’effet de la solution carbonatée et des maladies fongiques, chez le NAV TS par les dégâts des ravageurs, chez la laitue par les fongiques.
\begin{figure}[H]
    \centering
    \imgstd{images/v_longueur_chp5.png}
    \caption{Evolution moyenne de la longueur des feuilles en fonction du jour. }
    \label{fig:longueur-feuille5}
\end{figure}
Le gombo test porte des feuilles 2 à 3 fois plus longues que le témoin au jour 35 (9,8 cm contre 4,3 cm). Le navet test atteint 21,7 cm au J50, quand le témoin n'y arrive qu'au J67, soit 31~\% d'avance. L'allongement foliaire est plus favorable dans la partie test. L'effet ressort surtout sur le gombo et le navet.
\begin{figure}[H]
    \centering
    \imgstd{images/w_largeur_chap5.png}
    \caption{Evolution moyenne de la largeur des feuilles en fonction du jour. }
    \label{fig:largeur-feuille5}
\end{figure}
Les feuilles du gombo test sont plus épaisses que celles du témoin : 11,4~cm contre 6,3~cm à J35 (+81~\%), puis 14,0~cm contre 6,0~cm à J70 (+133~\%). Le navet test atteint 8,3~cm à J50, contre 7,0~cm pour le témoin à J67. La laitue test mesure 7,2~cm à J40. 
La largeur foliaire est plus développée dans la partie test, l'effet étant plus remarqué sur le gombo en fin de période.

{\small
\begin{longtable}{|p{2.5cm}|p{1.6cm}|p{6.1cm}|p{3.6cm}|}
    \caption{Synthèse des résultats des paramètres de croissance et de développement végétatif}
    \label{tab:synthese_croissance} \\
    \hline
    \textbf{Paramètre} & \textbf{Culture} & \textbf{Résultats principaux} & \textbf{Observation} \\
    \hline
    \endfirsthead
    \hline
    \textbf{Paramètre} & \textbf{Culture} & \textbf{Résultats principaux} & \textbf{Observation} \\
    \hline
    \endhead
    Survie des plants & Laitue & Le témoin disparaît à J20, tandis que le test conserve 5 plants à J40. & Survie prolongée dans le test. \\
    \hline
    Survie des plants & Gombo & Le témoin passe de 9 à 1 plant entre J14 et J27, soit environ 89~\% de mortalité. Le test conserve 4 plants de J50 à J70. & Forte mortalité dans le témoin. \\
    \hline
    Survie des plants & Navet & Le témoin passe de 6 à 3 plants entre J38 et J67. Le test reste stable. & Effectif plus stable dans le test. \\
    \hline
    Nombre de feuilles & Gombo & À J70 : 9,75 feuilles dans le test contre 5 dans le témoin. & Écart marqué en faveur du test. \\
    \hline
    Nombre de feuilles & Navet & 14,8 feuilles dans le test à J50 contre 13 dans le témoin à J67. & Valeur supérieure atteinte plus tôt dans le test. \\
    \hline
    Nombre de feuilles & Laitue & Le test atteint 9,8 feuilles à J40. & Développement foliaire observé dans le test. \\
    \hline
    Longueur des feuilles & Gombo & À J35 : 9,8~cm dans le test contre 4,3~cm dans le témoin. & Feuilles plus longues dans le test. \\
    \hline
    Longueur des feuilles & Navet & Le test atteint 21,7~cm à J50, valeur atteinte par le témoin à J67. & Développement plus précoce dans le test. \\
    \hline
    Largeur des feuilles & Gombo & À J35 : 11,4~cm contre 6,3~cm ; à J70 : 14~cm contre 6~cm. & Écart marqué en faveur du test. \\
    \hline
    Largeur des feuilles & Navet & 8,3~cm dans le test à J50 contre 7~cm dans le témoin à J67. & Valeur supérieure atteinte plus tôt dans le test. \\
    \hline
    Largeur des feuilles & Laitue & Le test atteint 7,2~cm à J40. & Développement foliaire observé dans le test. \\
    \hline
\end{longtable}
}

Dans l'ensemble, les observations montrent des différences de survie et de développement végétatif entre les plots tests et témoins. Les écarts les plus marques concernent le gombo, notamment le nombre et la largeur des feuilles. Le navet présente aussi un développement foliaire plus précoce dans le test. Pour la laitue, la mortalité précoce des plants témoins limite la comparaison des paramètres de croissance entre les deux modalités.

Ces résultats suggèrent une évolution plus favorable des cultures dans les plots TS. Toutefois, les différences de dates d'observations et d'effectifs de plants survivants doivent être prises en compte dans l'interprétation. Les attaques parasitaires observées sur le navet en fin de cycle constituent également un facteur susceptibles d'influencer son développent.

\begin{figure}[H]
    \centering
    \imgstd{images/x_hauteur_chap5.png}
    \caption{Evolution moyenne de la hauteur des feuilles en fonction du jour. }
    \label{fig:hauteur-feuille5}
\end{figure}
Le gombo test dépasse celui du témoin  environ +96,7~\% au jour 70 (41{,}3 contre 21{,}0~cm). L'écart apparaît tôt : dès le jour 35,  21{,}3~cm(gombo test) contre 7{,}0 cm(gombo témoin), soit le triple. Même chose pour le navet, le test la longueur de feuille atteint 16{,}9~cm à J41, contre 7{,}7~cm à J38 côté témoin, environ deux fois plus. La laitue test mesure 11{,}6~cm à J40, avant la mort de deux plants sur les sept marqués. La  partie test montre donc sa supériorité très tôt, et de façon marquée, sur le gombo et le navet en particulier.

\begin{figure}[H]
    \centering
    \imgstd{images/y_box_hauteur_chap5.png}
    \caption{Distribution de la hauteur des plantes selon la culture et le traitement. }
    \label{fig:hauteur-boxplot5}
\end{figure}
Les mesures ont montré que les plantes expérimentales de gombo et de navet ont des médianes de hauteur supérieures à celles des témoins. 
Cette différence est particulièrement marquante pour le gombo, avec une médiane de 28,0 cm contre 12,0 cm pour le témoin. 
Pour le navet, la médiane est également plus haute, à 22,0 cm contre 7,5 cm. Cependant, la hauteur du vaisseau testé a été interprétée comme la longueur des feuilles, ce qui rend cette comparaison moins équivalente. La laitue témoin n'est pas présente car elle n'a pas survécu.

\begin{figure}[H]
    \centering
    \imgstd{images/z_hauteur_final_chap5.png}
    \caption{Evolution de la hauteur moyenne des plants au dernier jour. }
    \label{fig:hauteur-jourfin}
\end{figure}
On évalue d'abord la hauteur maximale sur la parcelle de gombo test (41{,}2~cm), puis sur celle de gombo témoin (21{,}0~cm),suivie de la laitue test (11{,}6~cm), du navet test (10{,}4~cm) et, finalement, du navet témoin (9{,}2~cm).
Dans la partie test, la hauteur du gombo est deux fois plus grande que celle du  témoin et une stimulation  de la croissance du navet et de la laitue.

\begin{figure}[H]
    \centering
    \imgstd{images/z1_Effectif_hauteur_chap5.png}
    \caption{Distribution des effectifs par classe de hauteur. }
    \label{fig:effectif5}
\end{figure}
La bimodalité suggère une différenciation nette entre les deux types de sol. Les hauteurs faibles se regroupent dans la partie témoin, et les hauteurs élevées dans la partie test.

\begin{figure}[H]
    \centering
    \imgstd{images/z2_comparaison_L_l_chap5.png}
    \caption{Relation entre la longueur et la largeur de la plus grande feuille. }
    \label{fig:plus-G-feuille5}
\end{figure}
La figure ci-dessus montre la relation entre la longueur et la largeur de la plus grande feuille enregistrée sur chaque plante. 
Les points se regroupent en trois nuages distincts selon la culture. Le navet se trouve dans la zone des feuilles allongées, avec une longueur supérieure à 20 cm pour une largeur comprise entre 5 et 10 cm. 
Le gombo est dans la zone des feuilles larges, où la largeur atteint environ 17 cm pour des longueurs variant de 4 à 14 cm. La laitue, quant à elle, se situe dans la zone des feuilles plus petites, surtout en largeur. 
Pour tous les points combinés, le coefficient de corrélation de Pearson est de r = 0,31, ce qui indique une faible association positive. 
Ce résultat montre que la longueur et la largeur des feuilles sont faiblement liées de manière linéaire.

\begin{figure}[H]
    \centering
    \imgstd{images/z6_gom_nav_jour11_chap5.jpg}
    \caption{Photos comparatives des NAV T \& TS et GOM T \& TS au onzième jour. }
    \label{fig:gom-nav-j11}
\end{figure}

\begin{figure}[H]
    \centering
    \imgstd{images/z9_gombo_jour22_chap5.jpg}
    \caption{Photos comparatives des GOM T \& GOM TS au vingt-deuxième jour. }
    \label{fig:gombo-j22}
\end{figure}
Cette figure ci-dessus montre la différence entre le plot GOM T et GOM TS, les cercles avec contour rouge sur le GOM T indiquent les poquets qui contiennent toujours des cultures (les gombos qui ont résisté jusqu'à ce jour malgré les conditions difficiles)

\begin{figure}[H]
    \centering
    \imgstd{images/z10_navet_jour22_chap5.jpg}
    \caption{Photos comparatives des NAV T \& NAV TS au vingt-deuxième jour. }
    \label{fig:navet-j22}
\end{figure}

\begin{figure}[H]
    \centering
    \imgstd{images/z11_gombo_jour34_chap5.jpg}
    \caption{Photos comparatives des GOM T \& GOM TS au trente-quatrième jour. }
    \label{fig:gombo-j34}
\end{figure}
La figure ci-dessus montre qu'un seul plant a survécu dans le plot témoin et dans le TS, il y a des pertes par rapport aux jours précédents. Ce qui est lié aux brûlures de la solution carbonatée.

\begin{figure}[H]
    \centering
    \imgstd{images/z12_navet_jour34_chap5.jpg}
    \caption{Photos comparatives des NAV T \& NAV TS au trente-quatrième jour. }
    \label{fig:navet-j34}
\end{figure}

\begin{figure}[H]
    \centering
    \imgstd{images/z13_laitue_jour25_chap5.jpg}
    \caption{Photos comparatives des LAI T \& LAI TS au vingt-cinquième  jour. }
    \label{fig:laitue-j25}
\end{figure}
Cette image montre non seulement la présence de plants de laitue sur le plot test,  mais aussi indique que malgré l’apport important d’eau 
lié aux fortes pluies, il y a une parfaite infiltration complète sur le test, contrairement au témoin ou l’eau reste stagnante.



















